\section{引言：从 RNN 到 CNN}

本节课由 Chris Manning 主讲，介绍两种用于自然语言处理的神经网络技术：\textbf{卷积神经网络（Convolutional Neural Networks, CNN）}和\textbf{树递归神经网络（Tree Recursive Neural Networks）}。虽然在 Transformer 主导的当下，这两种技术的使用频率有所下降，但它们所蕴含的思想仍然具有重要的启发价值。

\begin{knowledgebox}{课程背景与定位}
Manning 指出，在任何科学领域中，不同的思想和技术总是交替出现的。即使某种方法当前不是主流，了解它们仍然有价值，因为人们经常以新的方式重新发明和组合这些思想。CNN 和树递归神经网络正是这样的例子——它们包含的核心思想至今仍然活跃在各种模型设计中。
\end{knowledgebox}

\subsection{RNN 的局限性}

回顾循环神经网络（RNN），它通过从左到右逐词处理来构建句子表示。这种方式有一个关键问题：

\begin{itemize}
    \item \textbf{无法捕获不带前缀上下文的短语}：RNN 在每一步都包含了之前所有词的信息。例如，对于句子``Monae walked into the ceremony''，你得到的不是``the ceremony''的表示，而是``Monae walked into the ceremony''的整体表示。
    \item \textbf{最后几个词的信息往往过度主导}：由于梯度消失等问题，最终的隐藏状态往往更多地反映句子末尾的内容。
\end{itemize}

\begin{figure}[H]
\centering
\includegraphics[width=0.95\textwidth]{slides-images/slide-005.jpg}
\caption{RNN 的局限：每一步的隐藏状态包含了所有前缀信息，无法单独表示局部短语\protect\footnotemark}
\end{figure}
\footnotetext{Slides 第6页。}

\subsection{CNN 的核心思想}

卷积神经网络提供了一种不同的思路：类似于 n-gram 模型，我们可以对固定长度的词子序列（如 bigram、trigram）计算向量表示，然后通过某种方式聚合这些表示。

\begin{figure}[H]
\centering
\includegraphics[width=0.95\textwidth]{slides-images/slide-006.jpg}
\caption{CNN 的基本思想：对每个固定长度的词子序列计算向量表示\protect\footnotemark}
\end{figure}
\footnotetext{Slides 第7页。}

\begin{importantbox}{CNN 用于 NLP 的核心思想}
对于句子``tentative deal reached to keep government open''，CNN 会计算所有 trigram 的向量表示：``tentative deal reached''、``deal reached to''、``reached to keep''等。这种方式不关心子序列是否是合法的语言学成分——它只是机械地对所有固定长度的窗口进行特征提取，然后通过后续操作进行聚合。
\end{importantbox}

\subsection{本章小结}

RNN 从左到右顺序处理文本，无法独立捕获局部短语的表示。CNN 通过滑动窗口的方式，为每个固定长度的子序列计算独立的向量表示，提供了一种互补的文本建模方案。

%% ========== 1D 卷积的基本操作 ==========

